\chapter{System Deployment and Testbench Architecture}
\label{ch:ch05_architecture_setup}

% --- Merged from ch08_forecasting_calibration.tex ---

\section{Forecasting and Calibration}
\label{ch:forecasting}

Dispatch safety begins with forecast quality: an optimiser can only be as
reliable as the prediction intervals it conditions on.  This chapter
presents the energy management instantiation of the ORIUS forecasting
layer.  The methods---six model families, conformal quantile regression,
OQE reliability scoring, and RAC interval inflation---are domain-agnostic;
their instantiation here uses the locked DE ENTSO-E load dataset as the
primary validation surface.  Chapter~\ref{ch:battery-to-orius} applies
the same forecasting layer to all six ORIUS domains at equivalent depth.

% ----------------------------------------------------------------
\subsection{Forecast Targets}

Three operational targets are forecast at hourly resolution over a
48-step rolling horizon ($H = 48$):

\begin{enumerate}
  \item \texttt{load\_mw} --- net system load (MW).
  \item \texttt{wind\_mw} --- aggregate wind generation (MW).
  \item \texttt{solar\_mw} --- aggregate PV generation (MW).
\end{enumerate}

A fourth target, \texttt{price\_eur\_mwh}, is forecast for the Germany
surface only and used in the dispatch objective
(Section~\ref{sec:dispatch-objective}, Chapter~\ref{ch:battery-dynamics}).
All targets are standardised to zero mean and unit variance before model
ingestion; predictions are de-standardised before interval construction.

% ----------------------------------------------------------------
\subsection{Model Families}
\label{sec:model-families}

Six model families span the spectrum from gradient-boosted trees to
attention-based deep architectures:

\begin{table}[ht]
\centering
\caption{Forecast model families and key hyper-parameters.}
\label{tab:model-families}
\begin{tabular}{llll}
\toprule
\textbf{Family} & \textbf{Type} & \textbf{Key hyper-params} &
  \textbf{Notes} \\
\midrule
GBM        & Gradient-boosted trees & depth, lr, rounds &
  Strongest on locked DE dataset \\
LSTM       & Recurrent (gated)      & hidden, layers, dropout &
  Baseline sequential model \\
TCN        & Temporal convolution   & kernel, dilation, channels &
  Parallelisable alternative to LSTM \\
N-BEATS    & Basis-expansion blocks & stacks, width, sharing &
  Interpretable decomposition \\
TFT        & Temporal fusion        & heads, hidden, quantiles &
  Multi-horizon with variable selection \\
PatchTST   & Patch-based transformer& patch len, heads, depth &
  Channel-independent patching \\
\bottomrule
\end{tabular}
\end{table}

\paragraph{GBM (LightGBM).}
Gradient Boosted Machines trained via LightGBM achieve the lowest MAE on
the locked Germany surface for load and wind targets.  The model ingests
94 engineered features including calendar dummies, rolling statistics, and
weather forecasts.  Quantile regression heads at $\tau \in
\{0.025, 0.05, 0.10, 0.25, 0.50, 0.75, 0.90, 0.95, 0.975\}$ are trained
independently, providing native prediction intervals before conformal
calibration.

\paragraph{LSTM.}
A two-layer LSTM with 128 hidden units and 10\% dropout serves as the
recurrent baseline.  Input sequences of 168\,h (one week) are consumed to
produce a 48-step ahead forecast.  The LSTM captures temporal dependencies
but is prone to over-smoothing during rapid ramp events.

\paragraph{TCN.}
A Temporal Convolutional Network with dilated causal convolutions
(kernel size~7, 4 dilation layers) provides a parallelisable alternative
to the LSTM.  TCN achieves comparable accuracy but trains 3--5$\times$
faster due to the absence of sequential hidden-state propagation.

\paragraph{N-BEATS.}
The Neural Basis Expansion Analysis architecture decomposes the forecast
into trend and seasonality stacks using learnable polynomial and Fourier
basis functions.  The interpretable variant is used, with 3 trend stacks
(polynomial degree~3) and 3 seasonality stacks (harmonics~5).

\paragraph{TFT.}
The Temporal Fusion Transformer combines variable selection, gated residual
connections, and multi-head attention for multi-horizon quantile forecasting.
TFT is the only architecture in the family that natively produces
probabilistic forecasts without a separate quantile-regression head.

\paragraph{PatchTST.}
The Patch Time-Series Transformer segments the input sequence into
non-overlapping patches of length~16, embeds each patch, and applies a
standard transformer encoder with 4 heads and 3 layers.
Channel-independent processing avoids cross-variable leakage and reduces
parameter count.

% ----------------------------------------------------------------
\subsection{Training Protocol and Leakage Control}
\label{sec:training-protocol}

Temporal leakage is the primary threat to forecast evaluation integrity.
The following protocol prevents it:

\begin{enumerate}
  \item \textbf{Expanding-window split.}  Training data grows monotonically;
    the test window is always strictly future.  No shuffle split is permitted.
  \item \textbf{Embargo gap.}  A 48-hour embargo separates the end of the
    training window from the start of the test window, ensuring that no
    feature derived from a lagged target can leak future information.
  \item \textbf{Walk-forward evaluation.}  The test set is divided into
    non-overlapping 168-hour (one-week) windows.  Metrics are computed per
    window and then aggregated (mean, P95).
  \item \textbf{Feature freeze.}  The feature engineering pipeline is frozen
    before model training begins.  No feature is constructed using information
    from the test period.
  \item \textbf{Seed locking.}  All random seeds (NumPy, PyTorch, LightGBM)
    are fixed at 42 for the primary run, with reproducibility verified across
    seeds $\{11, 22, 33, 44\}$.
\end{enumerate}

% ----------------------------------------------------------------
\subsection{Conformal Prediction Intervals}
\label{sec:conformal}

Point forecasts are converted to prediction intervals using Conformalized
Quantile Regression (CQR), following the split-conformal framework:

\begin{enumerate}
  \item Train quantile regressors at levels $\alpha/2$ and $1 - \alpha/2$
    on the training fold to obtain $\hat{q}_{\alpha/2}(x)$ and
    $\hat{q}_{1-\alpha/2}(x)$.
  \item On a held-out calibration fold of size~$n$, compute the conformity
    scores:
    \begin{equation}
    \label{eq:cqr-scores}
      s_i = \max\!\bigl(\hat{q}_{\alpha/2}(x_i) - y_i,\;
                        y_i - \hat{q}_{1-\alpha/2}(x_i)\bigr),
      \quad i = 1, \ldots, n.
    \end{equation}
  \item Let $\hat{q}$ be the $\lceil (1-\alpha)(n+1) \rceil / n$
    quantile of $\{s_1, \ldots, s_n\}$.
  \item The prediction interval for a new point~$x_{n+1}$ is:
    \begin{equation}
    \label{eq:cqr-interval}
      C(x_{n+1}) = \bigl[\hat{q}_{\alpha/2}(x_{n+1}) - \hat{q},\;
                         \hat{q}_{1-\alpha/2}(x_{n+1}) + \hat{q}\bigr].
    \end{equation}
\end{enumerate}

CQR provides a finite-sample marginal coverage guarantee:
$\Pr[y_{n+1} \in C(x_{n+1})] \ge 1 - \alpha$, assuming exchangeability
of the calibration and test data.

% ----------------------------------------------------------------
\subsection{Marginal Coverage}

On the locked Germany surface, the GBM + CQR pipeline achieves
PICP@90 $= 100\%$ and PICP@95 $= 100\%$ for load (marginal coverage).
Wind and solar targets also achieve marginal coverage above the nominal
level.  These aggregate numbers, however, mask conditional failures.

% ----------------------------------------------------------------
\subsection{Reliability-Conditioned Under-Coverage}

When coverage is stratified by the DC3S reliability score~$w_t$
(Chapter~\ref{ch:dc3s-adapter}), the picture changes substantially.
Table~\ref{tab:cqr-group-coverage} reproduces the locked group-coverage
results.

\begin{table}[ht]
\centering
\caption{CQR group coverage at 90\% nominal level, stratified by
reliability group.  Source:
\texttt{reports/publication/cqr\_group\_coverage.csv}.}
\label{tab:cqr-group-coverage}
\begin{tabular}{llrr}
\toprule
\textbf{Target} & \textbf{Group} & \textbf{PICP@90} &
  \textbf{Mean width (MW)} \\
\midrule
load\_mw  & low  & 0.802 & 2\,147 \\
load\_mw  & mid  & 0.825 & 2\,305 \\
load\_mw  & high & 0.807 & 2\,345 \\
\midrule
wind\_mw  & low  & 0.488 & 464 \\
wind\_mw  & mid  & 0.396 & 285 \\
wind\_mw  & high & 0.410 & 381 \\
\midrule
solar\_mw & low  & 0.903 & 796 \\
solar\_mw & mid  & 0.887 & 869 \\
solar\_mw & high & 0.816 & 853 \\
\bottomrule
\end{tabular}
\end{table}

Wind intervals are catastrophically under-covered in all three groups
(39--49\% vs.\ 90\% nominal), confirming that marginal guarantees do not
transfer to conditional subsets.  Load intervals under-cover by 8--20
percentage points; solar intervals are closest to nominal but still
under-cover in the high-variability group.

% ----------------------------------------------------------------
\subsection{Solar and Wind Calibration Weaknesses}
\label{sec:calibration-weaknesses}

Two structural factors drive the calibration gap:

\paragraph{Wind non-stationarity.}
Wind generation exhibits regime switches (storm fronts, lulls) that
violate the exchangeability assumption underlying CQR.  The calibration
fold, drawn from a single contiguous period, may not contain sufficient
representation of extreme regimes, leading to narrow intervals that
fail during the test period.

\paragraph{Solar bimodality.}
Solar generation has a deterministic diurnal envelope (sunrise--sunset)
overlaid with a stochastic cloud-cover component.  During nighttime hours,
the true generation is identically zero, making any nonzero interval
width wasteful.  During daytime, cloud transients produce heavy-tailed
residuals that require wide intervals.  A single set of conformity
scores cannot efficiently cover both regimes simultaneously.

% ----------------------------------------------------------------
\subsection{Why Calibration Alone Cannot Close the Gap}

The group-coverage results demonstrate a fundamental limitation:
\emph{post-hoc conformal calibration on a static calibration fold cannot
guarantee conditional coverage under distribution shift or telemetry
degradation.}

Three pathways to closing the gap are pursued in subsequent chapters:

\begin{enumerate}
  \item \textbf{Reliability-adaptive inflation} (DC3S, Chapter~\ref{ch:dc3s-adapter}):
    the interval width is inflated in real time as a function of the
    observation quality score~$w_t$, the drift flag~$d_t$, and the stale
    count~$s_t$, compensating for the conditional under-coverage.
  \item \textbf{Fault-tolerant invariant tube} (FTIT, Chapter~\ref{ch:dc3s-adapter}):
    a robust state tube propagates worst-case SOC bounds forward, providing
    a safety backstop that does not depend on interval calibration.
  \item \textbf{Online conformal update} (ACI baseline, Chapter~\ref{ch:cpsbench}):
    the conformal quantile is updated online using the most recent
    residuals, partially adapting to distribution shift but without
    observation-quality feedback.
\end{enumerate}

% ----------------------------------------------------------------
\subsection{Wind Coverage Gap: Diagnosis, Bound, and Partial Mitigation}
\label{sec:wind-gap}

Wind PICP@90 of 39--49\% is the most prominent calibration failure in this
thesis.  This section documents the root cause, the scope of the gap, and
the partial mitigation already in place.

\paragraph{Root cause: regime non-stationarity.}
Wind generation follows strongly seasonal and synoptic regimes.  The CQR
calibration set is a fixed temporal window; when the test set contains a
regime not represented in calibration (e.g.\ a prolonged high-pressure
anticyclone producing near-zero wind for days), the conformal quantile
$\hat{q}_\alpha$ is badly underfit.  The residuals on test are heavy-tailed
relative to the calibration distribution, so 50\% of test points fall
outside an interval sized for a stationary distribution.

\paragraph{Scope of the gap.}
Wind conditional under-coverage is a known limitation, not a hidden failure.
Table~\ref{tab:cqr-group-coverage} documents the 39--49\% values explicitly.
The DC3S shield compensates via two mechanisms: (1)~the OQE reliability
weight $w_t$ drops when the forecast residual exceeds expected magnitude,
triggering RUI inflation; (2)~Mondrian CQR stratification by reliability
group (detailed in Chapter~\ref{ch:conditional-coverage}) partially recovers
conditional coverage within each stratum, though marginal wind coverage
remains below nominal.

\paragraph{Quantitative bound.}
From \path{reports/publication/reliability_group_coverage_phase3.csv},
the wind stratum with $w_t \geq 0.8$ achieves PICP@90 = 72\%; the stratum
with $w_t < 0.5$ achieves 31\%.  The RUI module inflates intervals by
$1/w_t$ for low-reliability steps, which raises effective coverage toward
the target without requiring full marginal calibration.  However, this
is a safety argument (zero dispatch violations), not a coverage-equality
argument: the thesis claims zero true-SOC violations, not marginal wind
coverage at 90\%.  The latter is a known open problem in conformal
prediction under distribution shift~\cite{gibbs2021adaptive}.

\paragraph{Consequence for safety.}
Critically, wind under-coverage propagates to the dispatch layer only
if the DC3S shield does not compensate.  Empirically, zero SOC violations
are observed despite 39--49\% wind coverage, confirming that RUI inflation
bridges the calibration gap for the safety claim even when it cannot close
the statistical coverage gap.

% ----------------------------------------------------------------
\subsection{Chapter Summary}

This chapter described the six forecast model families (GBM, LSTM, TCN,
N-BEATS, TFT, PatchTST), the leakage-controlled training protocol, and the
CQR conformal calibration layer.  Marginal coverage meets or exceeds the
nominal level on all targets, but reliability-stratified analysis reveals
severe conditional under-coverage, particularly for wind (39--49\% vs.\ 90\%
nominal).  Solar calibration is weakened by diurnal bimodality, and wind by
regime non-stationarity.  These structural calibration failures motivate the
observation-aware adaptive inflation developed in the next chapter.


% --- Merged from ch10_cpsbench_battery_track.tex ---

\section{CPSBench Multi-Domain Evaluation Framework}
\label{ch:cpsbench}

Standard forecasting benchmarks evaluate prediction accuracy in isolation:
they score point forecasts or intervals against held-out actuals without
ever closing the loop through a physical plant.  The CPSBench Battery Track
closes that loop, coupling forecasts to a battery dispatch controller and
evaluating safety on the \emph{true} (unobserved) state trajectory under
replayable telemetry faults.

% ----------------------------------------------------------------
\subsection{Why Standard Evaluation Is Insufficient}

A forecast benchmark that reports MAE and PICP answers only half the
question.  Consider two controllers:

\begin{itemize}
  \item Controller~A achieves PICP@90 = 92\% with narrow intervals,
    maximising economic value---but during the 8\% of uncovered steps,
    the dispatched action drives the battery SOC below the safe floor,
    triggering a BMS trip that takes the asset offline for 30 minutes.
  \item Controller~B achieves PICP@90 = 96\% with wider intervals,
    sacrificing 3\% of revenue---but never triggers a BMS trip.
\end{itemize}

Standard metrics rank~A above~B on interval sharpness and close to~B on
coverage.  Yet in operational terms, A is unsafe and B is safe.  The
missing ingredient is \emph{closed-loop evaluation on the true plant
state}, which requires:
\begin{enumerate}
  \item A plant simulator with deterministic SOC dynamics.
  \item A fault injector that degrades telemetry between the plant and the
    controller.
  \item Separate logging of truth-state and observed-state trajectories.
  \item Safety metrics computed exclusively on the truth trajectory.
\end{enumerate}

CPSBench provides all four.

% ----------------------------------------------------------------
\subsection{Dual-Trajectory Evaluation}

The central architectural choice in CPSBench is the maintenance of two
parallel trajectories at every step~$t$:

\begin{description}
  \item[Truth trajectory $\{z_t^{\mathrm{true}}\}$.]
    The plant simulator advances the true SOC using the
    \emph{executed} action and the true load/renewables signal.  This
    trajectory is logged but never exposed to the controller.

  \item[Observed trajectory $\{\tilde{z}_t\}$.]
    The fault injector corrupts the truth signal before delivering it to
    the controller.  The controller makes all decisions based on this
    degraded view.
\end{description}

Safety violations are defined as events where the truth-trajectory SOC
exits the safe envelope $[E_{\min}, E_{\max}]$.  This is the
\emph{true-SOC violation} (as opposed to an observed-SOC violation, which
may be masked by sensor error).

The dual-trajectory architecture ensures that the benchmark answers the
correct question: ``did the battery actually violate its SOC limits?''
rather than ``did the controller \emph{think} it violated?''

% ----------------------------------------------------------------
\subsection{Truth-State and Observed-State Logging}

At each step the CPSBench runner logs the following record to the episode
trace (cf.\ \texttt{reports/publication/dc3s\_step\_trace.csv}):

\begin{table}[ht]
\centering
\caption{Per-step logging schema in CPSBench.}
\label{tab:cpsbench-log-schema}
\begin{tabular}{lll}
\toprule
\textbf{Field} & \textbf{Source} & \textbf{Trajectory} \\
\midrule
\texttt{load\_observed\_mw}    & Fault injector   & Observed \\
\texttt{renewables\_observed\_mw} & Fault injector & Observed \\
\texttt{renewables\_true\_mw}  & Plant simulator  & Truth \\
\texttt{reliability\_w}        & OQE              & Observed \\
\texttt{drift\_flag}           & Page--Hinkley    & Observed \\
\texttt{inflation}             & RAC-Cert         & Computed \\
\texttt{proposed\_charge\_mw}  & Optimiser        & Decision \\
\texttt{proposed\_discharge\_mw} & Optimiser      & Decision \\
\texttt{safe\_charge\_mw}      & DC3S shield      & Decision \\
\texttt{safe\_discharge\_mw}   & DC3S shield      & Decision \\
\texttt{intervened}            & Shield           & Decision \\
\texttt{soc\_after\_mwh}       & Plant simulator  & Truth \\
\texttt{guarantee\_passed}     & Certificate      & Computed \\
\bottomrule
\end{tabular}
\end{table}

Every field is deterministic for a given (scenario, seed, controller)
triple, enabling exact reproduction of any individual step.

% ----------------------------------------------------------------
\subsection{Fault Taxonomy and Replayable Schedules}

CPSBench implements six fault types, each controlled by a severity parameter
and applied by a deterministic injector seeded from the episode RNG:

\begin{table}[ht]
\centering
\caption{Fault taxonomy with severity parameterisation.}
\label{tab:fault-taxonomy-cpsbench}
\begin{tabular}{llll}
\toprule
\textbf{Fault} & \textbf{Parameter} & \textbf{Sweep levels} &
  \textbf{Effect on $\tilde{z}_t$} \\
\midrule
Dropout       & Rate $p$         & 0, 0.05, 0.10, 0.20, 0.30 &
  $\tilde{z}_t = \mathrm{NaN}$ w.p.\ $p$ \\
Delay jitter  & Max delay (s)    & 0, 1, 5, 15 &
  $\tilde{z}_t = z_{t-\delta}$, $\delta \sim U[0, d_{\max}]$ \\
Spike         & Scale $\sigma$   & 0, 1, 2, 3 &
  $\tilde{z}_t = z_t \cdot (1 + \epsilon)$, $\epsilon \sim N(0, \sigma^2)$ \\
Stale sensor  & Window (steps)   & 0, 2, 4, 8 &
  $\tilde{z}_t = z_{t-s}$ for $s$ consecutive steps \\
Drift         & Rate (MW/step)   & 0, 0.5, 1.0, 2.0 &
  $\tilde{z}_t = z_t + \beta \cdot t$ \\
Reorder       & Rate $r$         & 0, 0.05, 0.15 &
  Adjacent packets swapped w.p.\ $r$ \\
\bottomrule
\end{tabular}
\end{table}

Each fault schedule is a deterministic function of the episode seed,
meaning that any researcher with the same code and seed will reproduce
exactly the same fault pattern.  The \texttt{drift\_combo} scenario
applies all six faults simultaneously at moderate severity, representing a
worst-case but field-plausible operating condition.

% ----------------------------------------------------------------
\subsection{Benchmark Metrics}

The CPSBench Battery Track reports three families of metrics, computed
exclusively on the truth trajectory:

\subsubsection{Safety Metrics}

\begin{description}
  \item[TSVR (True-SOC Violation Rate).]
    Fraction of steps where the true SOC exits $[E_{\min}, E_{\max}]$:
    \[
      \mathrm{TSVR} = \frac{1}{T}\sum_{t=1}^{T}
        \mathbf{1}\!\bigl[E_t^{\mathrm{true}} < E_{\min}
          \;\lor\; E_t^{\mathrm{true}} > E_{\max}\bigr].
    \]
    The primary safety metric.  Any controller with $\mathrm{TSVR} > 0$
    has experienced at least one real SOC excursion.

  \item[Severity (P95).]
    The 95th percentile of violation magnitude across violating steps:
    \[
      \mathrm{Sev}_{95} = Q_{0.95}\!\Bigl(\bigl\{
        |E_t^{\mathrm{true}} - E_{\mathrm{nearest\_bound}}|
        : E_t^{\mathrm{true}} \notin [E_{\min}, E_{\max}]\bigr\}\Bigr).
    \]
    Measured in MWh, severity captures how far outside the safe envelope
    the battery ventured.

  \item[IR (Intervention Rate).]
    Fraction of steps where the shield modified the candidate action:
    \[
      \mathrm{IR} = \frac{1}{T}\sum_{t=1}^{T}
        \mathbf{1}\!\bigl[a_t^{\mathrm{safe}} \neq a_t^\star\bigr].
    \]
\end{description}

\subsubsection{Economic Metrics}

\begin{description}
  \item[Useful work (MWh).]
    Total energy throughput that generated non-negative revenue.
  \item[Expected cost (USD).]
    Dispatch cost under the tariff schedule, including degradation penalty.
  \item[Cost delta (\%).]
    Percentage change in cost relative to the deterministic-LP baseline.
\end{description}

\subsubsection{Environmental Metrics}

\begin{description}
  \item[Carbon (kg CO$_2$).]
    Grid carbon displaced by battery discharge, computed using the
    marginal emission factor from the generation-mix model.
\end{description}

% ----------------------------------------------------------------
\subsection{Benchmark Invariants}

The following invariants are checked programmatically after each episode
and recorded in the certificate stream:

\begin{enumerate}
  \item \textbf{Energy conservation.}  The truth-trajectory SOC at
    step~$t+1$ must equal the SOC at step~$t$ plus the energy balance
    from the executed action, to within floating-point tolerance
    ($\epsilon < 10^{-6}$\,MWh).
  \item \textbf{Power bounds.}  The executed power must satisfy
    $0 \le P_t^{+} \le P_{\max}^{+}$ and
    $0 \le P_t^{-} \le P_{\max}^{-}$ at every step.
  \item \textbf{Mutual exclusion.}  $P_t^{+} \cdot P_t^{-} = 0$ at every
    step (no simultaneous charge and discharge).
  \item \textbf{Certificate completeness.}  Every step must emit a
    certificate with zero missing fields (tracked as
    \texttt{certificate\_missing\_fields} in the results CSV).
  \item \textbf{Determinism.}  Re-running the same (scenario, seed,
    controller) triple must produce bit-identical results.
\end{enumerate}

% ----------------------------------------------------------------
\subsection{Battery Track as Benchmark Foundation}

The CPSBench Battery Track serves as the empirical foundation for all
subsequent evaluation chapters in this dissertation:

\begin{itemize}
  \item \textbf{Chapter~\ref{ch:cert-half-life-blackout}} uses the
    \texttt{drift\_combo} scenario to measure certificate half-life under
    communication blackout.
  \item \textbf{Chapter~\ref{ch:orius-bench-battery}} extends the battery
    track into a public benchmark release with standardised controller
    interfaces and reproduction instructions.
  \item The locked sweep results in
    \texttt{reports/publication/cpsbench\_merged\_sweep.csv} contain
    the full four-dimensional severity sweep (dropout rate $\times$
    delay $\times$ reorder rate $\times$ spike~$\sigma$) across all
    eight controllers and five seeds.
  \item The fault breakdown in
    \texttt{reports/publication/dc3s\_fault\_breakdown.csv} isolates
    the per-fault-type violation and intervention rates for each
    controller.
\end{itemize}

The battery track is intentionally \emph{single-asset}: it evaluates one
battery on one dataset surface at a time.  Fleet-level and cross-region
extensions are addressed in Chapters~\ref{ch:orius-bench-battery}
and~\ref{ch:cert-half-life-blackout}.

% ----------------------------------------------------------------
\subsection{Chapter Summary}

This chapter introduced the CPSBench Battery Track, a closed-loop
evaluation harness that couples forecasts to battery dispatch under
replayable telemetry faults.  The dual-trajectory architecture separates
truth state from observed state, enabling safety metrics (TSVR, severity)
to be computed on the actual plant trajectory rather than the controller's
degraded view.  Six fault types---dropout, delay jitter, spike, stale
sensor, drift, and reorder---are swept across configurable severity levels
with deterministic episode seeds.  Three metric families (safety, economic,
environmental) provide a comprehensive evaluation surface.  Benchmark
invariants (energy conservation, power bounds, mutual exclusion, certificate
completeness, determinism) ensure result integrity.  The battery track
forms the empirical foundation for all subsequent evaluation and extension
chapters.


% --- Merged from ch30_orius_bench_battery_track.tex ---
\section{ORIUS-Bench: Battery Track and Benchmark Contract}
\label{ch:orius-bench-battery}

This chapter promotes ORIUS-Bench from a mention of benchmark release into a
defended contract chapter.  The battery track matters not because it proves
universal superiority, but because it defines the reusable API, replay
protocol, metric schema, and artifact bundle that make observation-aware
safety evaluation reproducible.

\subsection{Benchmark Scope and Design Goals}

ORIUS-Bench (Battery Track) is an open, deterministic evaluation harness
designed to answer a single question: \emph{does a battery dispatch
controller remain safe when its telemetry is degraded?}

Design goals:
\begin{enumerate}
  \item \textbf{Reproducibility}: deterministic episode generation from a
    (scenario, seed) pair, eliminating run-to-run variance.
  \item \textbf{Fault isolation}: each fault dimension (dropout, delay,
    spike, stale sensor, drift) is independently controllable.
  \item \textbf{Truth separation}: the harness maintains parallel truth
    and observed trajectories, enabling unambiguous violation attribution.
  \item \textbf{Multi-controller comparison}: any controller conforming
    to the benchmark contract can be plugged in.
\end{enumerate}

\subsection{Episode Generation and Fault Replay}

Episodes are generated by the battery benchmark scripts and replayed under
named fault schedules.  The battery track carries independent severity knobs
for dropout, delay-jitter, out-of-order packets, spike magnitude, stale
sensors, and drift combinations.  This matters because the replay protocol is
part of the contribution: a score without a reproducible fault schedule is
not benchmark evidence.

\subsection{Seven-Metric Schema}

The ORIUS-Bench seven-metric schema
(\texttt{src/orius/orius\_bench/metrics\_engine.py}) exports:

\begin{enumerate}
  \item True-State Violation Rate (TSVR),
  \item Observation-Action Safety Gap (OASG),
  \item Certificate Validity Accuracy (CVA),
  \item Graceful Degradation Quality (GDQ),
  \item Intervention Rate (IR),
  \item Audit Completeness (AC), and
  \item Recovery Latency (RL).
\end{enumerate}

\subsection{Schema and Adapter Contract}

The benchmark exposes four schema groups through the locked artifact
\texttt{reports/publication/orius\_bench\_battery\_schema\_table.csv}.

\begin{table}[ht]
\centering
\caption{ORIUS-Bench battery schema groups and representative fields.}
\label{tab:orius-bench-schema}
\begin{tabular}{lll}
\toprule
\textbf{Schema group} & \textbf{Field} & \textbf{Meaning} \\
\midrule
Controller API & \texttt{controller} & controller identifier \\
Controller API & \texttt{safe\_action\_mw} & safety-repaired action power \\
Fault schema & \texttt{fault\_dimension} & injected fault family \\
Fault schema & \texttt{severity} & severity bucket \\
Metric schema & \texttt{true\_soc\_violation\_rate} & physical violation rate \\
Metric schema & \texttt{intervention\_rate} & safety intervention rate \\
Artifact schema & \texttt{run\_id} & execution lineage identifier \\
Artifact schema & \texttt{source\_script} & script that generated the run \\
\bottomrule
\end{tabular}
\end{table}

This is the reusable contribution of ORIUS-Bench: a domain adapter can vary
the plant and constraints, but it must still conform to the controller API,
fault replay contract, metric schema, and artifact schema.

\subsection{Controller Interface}

A benchmark-compliant controller must expose a \texttt{step()} method
returning at minimum:
\begin{itemize}
  \item \texttt{proposed\_charge\_mw}, \texttt{proposed\_discharge\_mw},
  \item \texttt{safe\_charge\_mw}, \texttt{safe\_discharge\_mw},
  \item \texttt{interval\_lower}, \texttt{interval\_upper}, and
  \item \texttt{certificate} or an explicit null certificate.
\end{itemize}

\subsection{Fault Replay and Determinism}

Battery leaderboard entries are only meaningful if replay is deterministic.
ORIUS-Bench therefore binds every run to:

\begin{enumerate}
  \item a scenario name,
  \item an explicit seed,
  \item a named source script,
  \item a run identifier, and
  \item a release-family manifest.
\end{enumerate}

The supporting governance artifacts live in
\texttt{reports/publication/run\_id\_policy.json} and
\texttt{reports/publication/release\_manifest.json}.  The benchmark result is
therefore not just a score table; it is a replayable artifact family.

\subsection{Leaderboard Surface}

Table~\ref{tab:orius-bench-leaderboard} reproduces the locked battery
leaderboard excerpt from
\texttt{reports/publication/orius\_bench\_battery\_leaderboard.csv}.

\begin{table}[ht]
\centering
\caption{ORIUS-Bench battery leaderboard excerpt.}
\label{tab:orius-bench-leaderboard}
\begin{tabular}{lrrrr}
\toprule
\textbf{Controller} & \textbf{Mean TSVR} & \textbf{Mean IR} &
  \textbf{Mean cost (USD)} & \textbf{Mean PICP@90} \\
\midrule
\texttt{scenario\_mpc} & 0.0000 & 0.0000 & 114346551.17 & 0.9347 \\
\texttt{scenario\_robust} & 0.0000 & 0.0000 & 114346551.17 & 0.9750 \\
\texttt{aci\_conformal} & 0.0000 & 0.0208 & 114365740.41 & 0.9653 \\
\texttt{dc3s\_wrapped} & 0.0000 & 0.0000 & 114401679.43 & 0.8806 \\
\texttt{dc3s\_ftit} & 0.0000 & 0.0882 & 114510699.19 & 0.9306 \\
\texttt{deterministic\_lp} & 0.1299 & 0.0000 & 114338186.20 & 0.8799 \\
\bottomrule
\end{tabular}
\end{table}

The leaderboard should be interpreted carefully.  It is useful benchmark
evidence, but it is not identical to the main thesis claim.  The battery
thesis still rests on the locked CPSBench battery evidence; ORIUS-Bench adds
release discipline, ranking surfaces, and replayability.

\subsection{Ranking Under Degradation}

The companion artifact \texttt{reports/publication/battery\_rank\_reversal.csv}
shows that degraded-observation ranking diverges from nominal ranking.  In
the current battery package, \texttt{dc3s\_wrapped} moves from nominal rank~7
to degraded rank~1, while \texttt{robust\_fixed\_interval} falls from nominal
rank~1 to degraded rank~7.  That rank reversal is precisely why the battery
track needs truth-vs-observation separation instead of nominal scorecards.

\subsection{What ORIUS-Bench Does and Does Not Prove}

ORIUS-Bench proves three things well:

\begin{enumerate}
  \item the battery evaluation protocol is reproducible,
  \item the controller and adapter contract is explicit, and
  \item degraded-observation ranking differs materially from nominal ranking.
\end{enumerate}

It does \emph{not} by itself prove cross-domain superiority.  The additional
adapters and unified harness are framework portability and auditability
evidence; the locked multi-domain theorem-and-artifact surface provides a
proof anchor for the thesis.

\subsection{Summary}

ORIUS-Bench (Battery Track) now appears in the thesis as a benchmark
contract, not just as a release note.  Its defended value is replayability,
schema discipline, and truth-separated evaluation rather than universal
dominance rhetoric.


% --- Merged from ch22_latency_systems_footprint.tex ---

\section{Runtime Latency and Systems Footprint}
\label{ch:latency}

Safety guarantees are worthless if the controller cannot compute them
within the dispatch cycle.  This chapter benchmarks the wall-clock
latency of every DC3S sub-component, characterises the runtime memory
footprint, and establishes that the full pipeline fits comfortably
within a 30-second dispatch interval on commodity hardware.

\subsection{Why Latency Matters for Real-Time Dispatch}

Grid-scale BESS dispatch operates on intervals of 5--30 minutes for
energy arbitrage and 2--4 seconds for frequency regulation.  The DC3S
pipeline must complete within the shortest applicable interval.  If the
safety shield exceeds the cycle time, the controller either dispatches
uncertified actions or stalls, both unacceptable outcomes.

\subsection{Measurement Methodology}

Latency is measured using Python \texttt{time.perf\_counter\_ns()} around
each DC3S sub-component during a 168-hour CPSBench episode under
\texttt{drift\_combo} (seed~42).  The episode comprises $N=168$ hourly
steps.  All measurements are recorded in
\texttt{reports/publication/dc3s\_latency\_summary.csv} and reproduced
below.

\subsection{Per-Component Latency}

Table~\ref{tab:latency} reports the wall-clock latency of each DC3S
pipeline stage.

\begin{table}[ht]
\centering
\caption{Per-component DC3S latency statistics (milliseconds).
  Source: \texttt{reports/publication/dc3s\_latency\_summary.csv}.}
\label{tab:latency}
\begin{tabular}{lrrrrr}
\toprule
\textbf{Component} & \textbf{Mean} & \textbf{Median} &
  \textbf{P95} & \textbf{P99} & \textbf{Max} \\
\midrule
Reliability scoring (OQE)     & 0.026 & 0.021 & 0.043 & 0.130 & 0.712 \\
Drift update (Page-Hinkley)   & 0.000 & 0.000 & 0.001 & 0.001 & 0.160 \\
Uncertainty set build (RUI)   & 0.012 & 0.011 & 0.011 & 0.047 & 0.724 \\
Action repair (SAF/Shield)    & 0.002 & 0.002 & 0.002 & 0.003 & 0.210 \\
\midrule
\textbf{Full DC3S step}       & \textbf{0.042} & \textbf{0.037} &
  \textbf{0.041} & \textbf{0.193} & \textbf{2.416} \\
\bottomrule
\end{tabular}
\end{table}

\subsection{Statistical Summary}

The full DC3S step completes in a \textbf{median of 0.037\,ms} and
a \textbf{P99 of 0.193\,ms}.  Even the worst-case single step
(2.416\,ms) is three orders of magnitude below the 30-second dispatch
interval.  The latency budget is dominated by reliability scoring
(62\% of mean time), followed by the uncertainty-set build (29\%).
The action repair projection is negligible ($<5\%$).

\subsection{Runtime Memory Footprint}

The DC3S pipeline maintains a sliding window of CQR residuals
(default: 200 points), the Page-Hinkley drift accumulator, and the
current certificate state.  Total resident memory is approximately
$\sim$12\,MB for a single-battery instance, measured via
\texttt{tracemalloc} snapshots:

\begin{itemize}
  \item CQR residual buffer: $\sim$3.2\,MB
    (200 points $\times$ 3 targets $\times$ quantile arrays).
  \item OQE feature history: $\sim$1.8\,MB.
  \item Forecast model (LightGBM): $\sim$5.6\,MB (pre-loaded).
  \item Certificate and state overhead: $\sim$1.4\,MB.
\end{itemize}

This footprint fits within the 64--256\,MB RAM available on typical
edge gateways (e.g., NVIDIA Jetson Nano, Raspberry Pi~4).

\subsection{Solver Infeasibility Rate}

The LP solver (HiGHS via \texttt{scipy.optimize.linprog}) reports
infeasible status when the tightened SOC bounds leave no feasible
dispatch.  Across the full CPSBench sweep, the solver-OK rate for
DC3S+FTIT is $\ge$95.8\%, with infeasibilities occurring only during
extreme drift episodes where the uncertainty set contracts the
feasible region to near-zero width.  In these cases, the SAF module
falls back to a safe-standby action (zero net power), preserving
safety at the cost of one lost dispatch cycle.

\subsection{Observability Outputs}

Each DC3S step emits a structured certificate containing:
\begin{itemize}
  \item Proposed and safe actions (charge/discharge MW).
  \item Reliability score $w_t$ and its component features.
  \item CQR interval width and adaptive multiplier.
  \item Solver status and objective value.
  \item SHA-256 certificate hash for audit traceability.
\end{itemize}

These outputs are logged to \texttt{reports/hil/hil\_certificate\_audit.csv}
(Chapter~\ref{ch:hil}) and can be streamed to an observability backend
(e.g., Prometheus/Grafana) for real-time monitoring.

\subsection{Deployment Implications}

The sub-millisecond median latency has three practical consequences:

\begin{enumerate}
  \item \textbf{Frequency regulation}: DC3S can serve 4-second AGC
    signals with $>$99.99\% on-time delivery.
  \item \textbf{Fleet scaling}: At 0.042\,ms per battery, a single
    core can certify $\sim$700 batteries per 30-second cycle, enabling
    fleet-level deployment without parallelisation
    (cf.\ Chapter~\ref{ch:compositional-safety-fleets}).
  \item \textbf{Edge deployment}: The combined latency and memory
    footprint are compatible with ARM-based edge gateways, eliminating
    cloud round-trip dependencies.
\end{enumerate}

\subsection{Summary}

The DC3S safety pipeline completes in a median of 0.037\,ms per step,
with worst-case latency of 2.4\,ms.  The full pipeline including
forecast, OQE, RUI, SAF, and certificate generation fits within
$\sim$12\,MB RAM.  These figures confirm that the safety guarantee
is computationally free in the context of grid-scale dispatch cycles,
removing latency as a barrier to deployment.


